% =====================================================================
%  CATALOGO DELLA BANCA ESERCIZI
%  Compila questo file per vedere tutti gli esercizi con soluzioni,
%  nome del file e controllo dei risultati delle file A, B e C.
%  Quando aggiungi un esercizio alla banca, aggiungi qui una riga.
% =====================================================================
\def\versione{completa}
\providecommand{\fila}{A}       % metti B o C per vedere i dati delle altre file

\documentclass[a4paper,11pt,addpoints]{exam}
\usepackage{stili/verifica}
\tipoprova{Catalogo esercizi}

\begin{document}
\begin{center}{\LARGE\bfseries Catalogo esercizi\par}\end{center}
\bigskip

\begin{questions}

% ---------- Elettrotecnica -------------------------------------------
\fileesercizio{Banca_Esercizi/elettrotecnica/linea_rame_01}
\fileesercizio{Banca_Esercizi/elettrotecnica/serie_incognita_01}
\fileesercizio{Banca_Esercizi/elettrotecnica/parallelo_incognita_01}
\fileesercizio{Banca_Esercizi/elettrotecnica/mista_E_da_I2_01}
\fileesercizio{Banca_Esercizi/elettrotecnica/scala_E_da_I5_01}
\fileesercizio{Banca_Esercizi/elettrotecnica/linea_parallelo_E_da_I2_01}
\fileesercizio{Banca_Esercizi/elettrotecnica/due_rami_E_da_I4_01}
\fileesercizio{Banca_Esercizi/elettrotecnica/mista_R_da_I_01}
\fileesercizio{Banca_Esercizi/elettrotecnica/mista_doppio_parallelo_E_da_I4_01}
\fileesercizio{Banca_Esercizi/elettrotecnica/kirchhoff_risalita_01}
\fileesercizio{Banca_Esercizi/elettrotecnica/ohm_potenza_01}
\fileesercizio{Banca_Esercizi/elettrotecnica/parallelo_tre_01}
\fileesercizio{Banca_Esercizi/elettrotecnica/circuito_misto_01}

\end{questions}

\clearpage
\tabellacontrollo
\end{document}
